\chapter{Limits and known problems}
\label{ch:limits}

This chapter lists what BAG does not do, and what you must know to use its output
safely. Every item was checked against the code of this version. For each limit you find
what happens, what you can do about it, and the check rule ID or warning that reports it,
if there is one. A limit with no rule ID is not reported by \cmd{bag check}: you must
keep to it yourself.

\section{Protocols and bursts}

\begin{itemize}
\item \textbf{Protocol versions.} BAG knows five protocols: AXI4, AXI4-Lite, AHB-Lite,
  APB4 and APB3. There is no AXI3, AXI5, AHB5, APB5, full AHB (with SPLIT and RETRY) or
  AXI-Stream. AXI4 ports have no ATOP signal, so AXI5 atomic transactions are not
  supported. The master-port guards drive \code{AWATOP} to zero.
\item \textbf{WRAP bursts through a width converter.} The vendored PULP width converters
  cannot split a WRAP burst. A \code{dw} adapter answers every WRAP burst with SLVERR, and
  so does \code{bag_axi2ahb} when its two sides have different data widths. When the
  converter makes the bus narrower, a FIXED burst of more than one beat also gets
  SLVERR. \cmd{bag check} warns once for each AXI4 master and converter, with
  \ruleid{AXI-WRAP-REFUSED}. The message names the slaves and bridges that the master can
  reach through that converter.

  What you can do: send INCR bursts on that path, or keep the path at one data width and
  one protocol. The warning lists every target with a region on the master's fabric. It
  does not look at the master's address width, so it can name a target that the master
  cannot address.
\item \textbf{WRAP bursts into AXI4-Lite or APB.} The protocol converters from AXI4 to
  AXI4-Lite (a \code{pc} adapter, and the bridges \code{bag_axi2axil} and
  \code{bag_axi2apb}) also answer every WRAP burst with SLVERR. \cmd{bag check} does not
  warn here, because register buses rarely see WRAP bursts. The generated tests expect the
  SLVERR on these paths.
\item \textbf{Non-modifiable bursts.} A burst with \code{AxCACHE[1]} = 0 is
  non-modifiable: a converter may not split it. The AXI4 to AXI4-Lite converters (the
  \code{pc} adapter, \code{bag_axi2axil}, \code{bag_axi2apb}) answer with SLVERR a
  non-modifiable FIXED burst of more than one beat, and a non-modifiable INCR burst of 2 to
  17 beats. The two bridges also refuse a non-modifiable access that is wider than the
  downstream bus. \code{bag_axi2ahb} refuses a non-modifiable burst that its width
  converter must split. \cmd{bag check} warns where CACHE is tied to 0 on the way
  (\ruleid{AXI-SPLIT-REFUSED}): every access is then non-modifiable.

  What you can do: send modifiable bursts (\code{AxCACHE[1]} = 1) on these paths, or
  single beats.
\item \textbf{Write errors through a protocol converter.} The AXI4 to AXI4-Lite
  converters report every write error as SLVERR. A DECERR from behind them reaches the
  master as SLVERR.
\item \textbf{Arbitration.} The AXI4 and AXI4-Lite crossbars arbitrate round-robin only
  (\ruleid{FAB-ARBITRATION}). Shared buses and the AHB-Lite bus also offer fixed priority.
  No fabric arbitrates by QoS. An AHB-Lite fabric is a single-layer bus; there is no
  multi-layer matrix (\ruleid{FAB-TOPOLOGY}). An APB fabric has one master
  (\ruleid{FAB-SINGLE-MASTER}).
\item \textbf{ID widths.} Each master port of an AXI4 fabric has a guard that counts
  transactions per ID. Its tables double with each master-side ID bit.
  \ruleid{FAB-GUARD-ID-WIDTH} warns above 8 bits. Each crossbar level adds bits to the
  slave-side ID; \ruleid{FAB-ID-GROWTH} warns above 16 bits. Both are area costs, not
  errors. Narrow the master IDs if the area matters.
\item \textbf{Max mst trans and Max slv trans mean different things per protocol.} The
  vendored PULP crossbars count differently for AXI4 and AXI4-Lite, and BAG passes the
  values through. Table~\ref{tab:limits-trans} shows the meaning. No check reports a
  difference, because none is an error.
\end{itemize}

\begin{table}[htbp]
\centering\small
\caption[What Max mst trans and Max slv trans limit]{What \col{Max mst trans} and
  \col{Max slv trans} limit}
\label{tab:limits-trans}
\begin{tabularx}{\linewidth}{@{}>{\raggedright\arraybackslash}p{0.2\linewidth}
  >{\raggedright\arraybackslash}X >{\raggedright\arraybackslash}X@{}}
\toprule
Fabric & \col{Max mst trans} & \col{Max slv trans} \\
\midrule
AXI4 crossbar & At least this many per ID at each master, not in total. &
  Writes per slave whose data has not finished. Reads are not limited. \\
AXI4 shared bus & At least this many per ID at each master. Also the depth of the
  write FIFO of the bus, shared by all masters. & At least this many per ID, across all
  slaves. \\
AXI4-Lite crossbar or shared bus & This many per master. & This many per slave, reads
  and writes. \\
\bottomrule
\end{tabularx}
\end{table}

Both columns also size the adapters on the ports of the fabric. \col{Max slv trans} also
limits the reads and writes through a \code{bag_axi2apb} or \code{bag_axi2axil} bridge on
the fabric, and the reads through \code{bag_axi2ahb}.

\section{Bridges and adapters}

\begin{itemize}
\item \textbf{\code{bag_axi2ahb} handles one transaction at a time.} It serializes the
  IDs, so responses come back in order and the bridge never overlaps two transactions.
  A read stays open until all its beats are in a 4-entry read buffer.

  This has one rule for you. A master must not make RREADY wait for the progress of a
  later write (its AW, W or B). If it does, the bridge and the master wait for each other,
  and the bus stops.
\item \textbf{\code{bag_axi2ahb} and AHB-Lite features.} The bridge answers an exclusive
  access (\code{AxLOCK}) with OKAY, so the exclusive access fails. It drives
  \code{HMASTLOCK} to zero. It does not carry USER signals (they read as zero). An INCR
  burst becomes an AHB-Lite INCR burst of undefined length, which restarts at each 1 KB
  boundary. WRAP and FIXED bursts become single transfers. After an error response, the
  burst continues, and the master gets SLVERR.
\item \textbf{Partial writes into AHB-Lite become byte writes.} AHB-Lite has no write
  strobes. \code{bag_axi2ahb} and \code{bag_apb2ahb} split a write with some strobes
  clear into one byte write for each set strobe. Contiguous strobes, such as a halfword,
  are split too. \code{bag_axi2ahb} drops a write beat with no strobe set.

  So such a write is not atomic at the slave: the slave sees several writes, and another
  master can access the word between them. What you can do: write whole words to AHB-Lite
  registers where this matters. A change to this behaviour is planned.
\item \textbf{The bridges from AHB-Lite and APB do one transfer at a time.}
  \code{bag_ahb2apb}, \code{bag_ahb2axi}, \code{bag_apb2ahb} and \code{bag_apb2axi} pass
  each transfer through a four-phase handshake and wait for its response. Their
  \col{FIFO depth} has no effect, and \ruleid{BRIDGE-FIFO-IGNORED} warns if you set one.
  \code{bag_ahb2axi} and \code{bag_apb2axi} issue AXI ID width 1 with one transfer
  outstanding. A 64-bit AHB-Lite beat into a 32-bit AXI bus becomes two AXI transfers.
\item \textbf{Signals the PULP-based bridges do not carry.} \code{bag_axi2axil} and
  \code{bag_axi2apb} do not carry USER, \code{AxCACHE}, QoS, REGION or LOCK. An exclusive
  access through them reaches the slave as a normal access and gets OKAY, so it fails.
  \code{bag_axi2apb} does one APB transfer at a time.
\item \textbf{The AHB-Lite bus does not forward \code{HMASTLOCK}.} A locked sequence
  keeps the grant of the bus, but the slaves never see \code{HMASTLOCK}.
\item \textbf{Partial writes into APB3.} APB3 has no PSTRB. A partial write reaches an
  APB3 slave that is wider than 8 bits as a full-word write: the bytes with a clear
  strobe are written too. The same happens on an APB4 slave whose PSTRB you tie in
  \sheet{Port_Options}, and \cmd{bag check} warns about that tie (\ruleid{PORTOPT-PSTRB-WORD}).
  What you can do: write whole words to such slaves.
\item \textbf{Narrow APB slaves.} \code{bag_axi2apb} splits a beat that is wider than
  the APB bus into several APB transfers. An error on one of them does not stop the
  others. So a write can change some bytes of the slave while the master gets SLVERR.

  An unaligned access from an AXI4-Lite master behaves differently on an 8-bit and a
  16-bit APB bus. Onto 8 bits, the bridge accesses the bytes from the address to the end
  of the beat, and not the bytes below the address. Onto 16 bits, it accesses the whole
  aligned word. Both are legal AXI. What you can do: use aligned accesses to narrow APB
  slaves.
\item \textbf{AXI4-Lite width converters access the whole wider word.} An AXI4-Lite
  \code{dw} adapter always reads the whole wider word. It writes the whole wider word too,
  with \code{WSTRB} = 0 on the bytes the master did not write. So a read also reads
  bytes at the slave that the master did not address. Keep registers with read side effects (a FIFO pop, a clear-on-read
  status) away from the other half of the wider word.
\item \textbf{FIFOs and register slices exist only on AXI ports.} \sheet{Port_Options}
  accepts a \col{FIFO depth} or a \col{Register slice} only on an AXI4 or AXI4-Lite port.
  \ruleid{PORT-ADAPTER} reports any other.
\end{itemize}

\section{Address decode and aliases}

\begin{itemize}
\item \textbf{Alias-enable pins must be quasi-static.} An alias enable that is a
  top-level input goes straight into the decoder of the fabric. There is no synchronizer.
  If the pin changes while a transfer is in flight, BAG does not define which map
  decodes it. The generated tests and the formal jobs assume that the pin changes only
  while the fabric is in reset.

  What you can do: change an alias pin only in reset, or while no master of the fabric
  has a transfer open. If the pin comes from another clock domain, synchronize it to the
  fabric's clock yourself. \ruleid{DEC-ALIAS-ENABLE} checks only that the pin name is
  \code{always} or an extra top-level input.
\item \textbf{The map is static.} BAG writes the memory map and the access rules into the
  RTL. Software cannot change regions or permissions at run time, and there is no
  firewall with lock bits.
\item \textbf{What the access control covers.} BAG enforces \sheet{Access_Rules}, the
  read-only and write-only regions, and secure-only regions. The secure check reads only
  the non-secure bit of PROT (\code{AxPROT[1]}, \code{HNONSEC}, \code{PPROT[1]}). A
  secure-only region must not be reachable from a port whose PROT is tied, or whose
  protocol has no PROT: \ruleid{SEC-PROT-TIED} reports it.
\item \textbf{The default slave is yours.} BAG does not know what a default slave
  answers. Nothing warns when it answers a miss with OKAY. A tied PSLVERR on the master
  side of an APB fabric hides every error, its own miss response included
  (\ruleid{PORTOPT-ERROR-MASKED}).
\item \textbf{What a slave sees.} A slave sees the address after alias translation and
  after the offset of a bridge with \col{Pass-through address} N. A slave with a narrow
  address port sees only the low bits. Its regions must lie in one aligned block of its
  address width (\ruleid{MM-SLAVE-ALIGN}).
\end{itemize}

\section{Clocks, resets and clock-domain crossings}

\begin{itemize}
\item \textbf{The CDC constraints are not verified by a tool.} \cmd{bag generate} writes
  \file{constraints/<project>_cdc.sdc} and \file{.xdc} (section~\ref{sec:constraints}), but
  BAG runs no synthesis or timing tool, so nothing proves that a tool accepts them or that
  their names match its netlist. What remains yours:
  \begin{itemize}
  \item the files assume the fabric clocks are ports of the top wrapper
    (\cmd{create_clock} is on \cmd{get_ports}); if your SoC derives them, constrain them in
    your own file and drop these lines;
  \item the register patterns assume the \code{_reg} suffix that synthesis gives a
    register, and the pin patterns hold only while synthesis keeps the \code{bag_sync} and
    \code{axi_cdc} boundaries; PULP asks for \cmd{set_ungroup false} for the same reason;
  \item the paths from a reset port into \code{bag_rst_sync} are not constrained;
  \item there is no hook for a technology synchronizer cell.
  \end{itemize}
  The synchronizer flops of \code{bag_sync} and \code{bag_rst_sync} carry
  \code{ASYNC_REG} and \code{DONT_TOUCH}, one spelling for every tool; a tool that does not
  know them ignores them.
\item \textbf{Resets of a reset group assert together.} Domains joined by a chain of
  crossings form a reset group. All resets of a group must assert at the same moment.
  Resetting one domain of a group alone is not supported: the PULP CDC FIFOs need both
  sides in reset together. If they are not, the running side can see a transfer that no
  master sent. \ruleid{CLK-CDC-RESET} checks that each reset on a crossing has async
  assertion and sync deassertion. It cannot check when your resets assert.

  \cmd{bag generate} writes \file{doc/reset.md} (with \setting{Generate memory-map doc}
  Y). It lists the reset groups of your design and the reset sequence to follow. BAG
  sees only its own crossings: a crossing inside your user-supplied fabric is yours.
\item \textbf{No timeouts.} No fabric or bridge has a timeout. A slave that never raises
  PREADY, HREADYOUT or an AXI READY, or never sends its response, stops the master. On a
  shared bus it stops every master. What you can do: use slaves that always answer, and
  add a watchdog outside BAG where you need one. A PREADY timeout is planned.
\item \textbf{No clock-domain crossing analysis.} BAG runs no CDC tool. Simulation has
  zero delay, so it cannot show metastability or a missing constraint.
\end{itemize}

\section{Verification}

\begin{itemize}
\item \textbf{Simulation coverage is a report only.} \cmd{bag sim} writes
  \file{sim/coverage.md} (\file{sim_ready_shim/coverage.md} with \opt{--ready-shim}). It has no closure target. Its bins come from the same addresses
  and burst shapes that the tests use. So a full passing run reads 100\% by construction:
  the report says what a run exercised, not what the tests leave out. It has no bins for
  back-pressure per channel, arbitration contention, alias pins on and off, narrow single
  transfers, or narrow bursts outside AXI4.
\item \textbf{What the generated tests do not try.}
  \begin{itemize}
  \item WRAP and FIXED bursts are tested only at full width.
  \item The burst at a 4 KB edge is not run on a region of exactly 4 KB.
  \item Every reset is asserted together and released together. A staggered release,
    and assertion with stopped clocks, are not tested.
  \item \cmd{bag sim} uses one fixed seed (\code{20261031}) unless you set
    \code{BAG_COCOTB_SEED}. A run tries one seed only: the tests do not sweep seeds.
  \end{itemize}
\item \textbf{\cmd{bag sim} does not judge the build log.} It builds the design with
  Verilator and runs the tests, but it does not read the warnings of the build. An
  elaboration \code{$error} in the RTL shows in \cmd{bag lint}. \cmd{bag flow} runs lint
  before sim, so use it, or run \cmd{bag lint} first.
\item \textbf{Formal is bounded, except two jobs.} The decode job is a complete proof,
  because the decoder is combinational. The reset job is proved by induction, so it holds
  in every reachable state. The protocol job and the route jobs check only to a fixed
  number of solver steps (16 for the protocol job). A bug that needs more steps to show
  is not found. There is no liveness check: no job proves that every request gets a
  response.
\item \textbf{What the route jobs compare.} For each request at a slave port, a route
  job checks the address, the direction, the secure bit and the slave port. It does not
  check the other request fields (length, size, burst type, cache and protection bits),
  the write data, or the payload and routing of responses. On a fabric with several
  masters and no IDs (an AXI4-Lite crossbar or shared bus, or the AHB-Lite bus), the job does
  not explore two masters with open requests to the same address at the same time. The
  window check of a bridge uses single-beat accesses, one at a time.
\item \textbf{Clocks in formal.} The formal jobs treat every clock as a free input, so
  they explore any clock ratio and phase. They model a crossing as ideal logic: they prove
  nothing about metastability or the timing of a synchronizer.
\item \textbf{The PULP blocks are not proved inside.} The properties sit at the ports of
  the design, of each fabric and of each bridge. Nothing is proved about the internals of
  the vendored PULP blocks, and BAG does not change their code. Adapters have no route job
  of their own; the protocol job and the simulation see them only from the ports.
\item \textbf{User-supplied fabrics get no tests and no formal.} BAG does not know how a
  \code{bag_external} fabric behaves. It writes a stub with all outputs tied to zero, so
  the design lints and elaborates. It writes no tests and no formal jobs for such a
  design. The memory map, the decode tables and the documents are still generated. If
  fabrics share an \col{External module}, their ports must be equal;
  \ruleid{EXT-MODULE-SHARED} checks the protocol, widths and attachments, and
  \cmd{bag generate} stops on a difference that the check cannot see, such as an ID width.
\end{itemize}

\section{Tools}

\begin{itemize}
\item \textbf{Pinned versions.} BAG is tested with the tool versions in
  \file{tools.lock}: Verilator, slang, Verible, and SymbiYosys and Yosys from one OSS CAD
  Suite release. With another version, findings can differ. The built-in lint waivers
  were written for the pinned versions: with another version, a waiver that matches
  nothing is a note, not an error. \cmd{bag formal} notes a SymbiYosys that is not the
  pinned one.
\item \textbf{The EDA flows run in the tools image.} They are tested there, on
  \code{linux/amd64} and \code{linux/arm64} (Ubuntu 24.04). Reading, checking,
  generating and the collateral need no EDA tools and run on Linux, macOS and Windows.
\item \textbf{No native EDA path on macOS.} The pinned Verilator publishes no macOS
  binary and cannot be built from source on macOS~27 (section~\ref{sec:macos}), so the
  flows are supported there only through the image, which is native on Apple Silicon.
  \cmd{bag lint} and \cmd{bag formal} do in fact work with natively installed tools,
  since slang, Verible and the OSS CAD Suite all publish macOS builds; that is not a
  tested configuration and is not supported. Note also that without the pinned Verilator
  on \code{PATH}, a native \cmd{bag lint} silently uses the OSS CAD Suite's own
  Verilator, which is not the pinned one.
\item \textbf{What is untested on arm64.} The full list (\code{eda and not deep})
  passes on \code{linux/arm64}. The \emph{deep} formal proofs and \cmd{bag regress}
  have not been run there.
\item \textbf{One simulator.} \cmd{bag sim} runs Verilator with cocotb only.
  \setting{Simulator} accepts only \code{verilator} (\ruleid{VALUE-ILLEGAL}). Xcelium and
  VCS are not supported, and the bus models are not tested with them. \cmd{bag lint} runs
  slang, Verilator and Verible, and \cmd{bag formal} runs SymbiYosys. There are no
  commercial backends.
\item \textbf{Generated tests need the same BAG version.} The generated test module
  imports the tests from the installed BAG. If that is another version, the tests stop
  with a message that tells you to run \cmd{bag generate} again.
\end{itemize}

\section{Generated output}
\label{sec:limits-output}

\begin{itemize}
\item \textbf{Two subsystems need different names that you choose.} The module prefix does
  not rename \setting{Project name}, \setting{Top module name} or \setting{Parameter
  package}, nor the generated interfaces and flat harness that carry them. Two subsystems
  with the same ones declare the same units, and BAG cannot check that across runs
  (section~\ref{sec:two-subsystems}). The joint compile of two prefixed subsystems and an
  unmodified PULP copy is tested with slang and Verilator, against the same PULP version
  only.
\item \textbf{The environment's macros are shared by design.} The module prefix reaches
  every macro that a packaged file defines, but not the macros that only the environment
  defines and the files merely test: \code{VERILATOR}, \code{SYNTHESIS}, \code{UVM},
  \code{QUESTA}, \code{VCS}, \code{XSIM}, \code{TARGET_GENUS}, \code{FPV_ON},
  \code{ASSERTS_OFF}, \code{ASSERTS_OVERRIDE_ON} and \code{COMMON_CELLS_ASSERTS_OFF}.
  The vendored assertion header declares the package \code{assert_rpt_pkg} inside
  \code{`ifdef UVM}, and it is renamed like every other unit, so an SoC with a UVM
  environment sees one package per subsystem. No BAG flow defines \code{UVM}, so BAG's own
  flows never elaborate the renamed package.
\item \textbf{Sub-wrappers have flat ports.} A sub-wrapper from the \sheet{Hierarchy}
  sheet has flat ports, also when \setting{Port style} is \code{sv-interface}. It receives
  the conditioned reset of each domain, not the raw reset.
\item \textbf{Some columns are kept but not used.} BAG checks
  Masters > \col{Max outstanding rd} and \col{Max outstanding wr}, Slaves >
  \col{Wait states} and \col{Error response}, and Fabrics > \col{Backing IP}, and keeps
  them in the IR, but no generated file depends on them yet.
\end{itemize}

\section{Not in this version}

\begin{itemize}
\item \textbf{No functional-safety mechanisms.} BAG generates no bus timeouts, no
  parity or ECC on bus signals, no fabric duplication with a comparator, and no
  error-status registers. It gives no safety manual and no FMEDA data.
\item \textbf{No tool qualification.} There is no qualification kit and no
  tool-confidence evidence for a functional-safety standard. Deterministic output, the
  generated tests and the formal jobs are evidence you can use, but BAG does not package
  them for qualification.
\item \textbf{No plug-in templates.} You can use only the built-in fabric, adapter and
  bridge templates (\ruleid{TMPL-FABRIC}, \ruleid{TMPL-BRIDGE}). For a fabric that BAG
  cannot generate, use \code{bag_external} and supply the RTL. Plug-in templates are
  planned.
\item \textbf{Other things BAG does not do.} BAG does not import or export IP-XACT. It
  does not report area or timing. It does not support custom protocols.
\end{itemize}
